% !TeX program = xelatex
% SIGMOD 2027 research submission: standard ACM sigconf, author-identifiable working copy.
% Keep the class's text size, margins, column separation and line spacing.
\PassOptionsToPackage{safe}{pbalance}
\documentclass[sigconf,pbalance]{acmart}
\setcopyright{none}
\settopmatter{printacmref=false,printccs=false,printfolios=true,authorsperrow=2}
\acmConference[SIGMOD 2027 Submission]{SIGMOD 2027 Submission}{}{}
\acmYear{2027}
\acmDOI{}
\acmISBN{}
\usepackage{placeins}
\usepackage{tabularx}
\usepackage{array}
\newlength{\pairedfigureheight}
\usepackage[ruled,vlined,linesnumbered,noend]{algorithm2e}
\SetKwInput{KwIn}{Input}
\SetKwInput{KwOut}{Output}
\SetKwInput{KwAssume}{Assumption}
\SetKw{Continue}{continue}
\SetAlFnt{\footnotesize}
\SetAlCapFnt{\small}
\SetAlCapNameFnt{\small\bfseries}
\SetAlgoNlRelativeSize{-1}
\SetInd{0.4em}{0.7em}
\SetCommentSty{textit}
\newcommand{\AFn}[1]{\operatorname{\mathit{#1}}}
\newcommand{\AState}[1]{\mathsf{#1}}
% Allow an extra line-breaking pass for narrow columns; no layout dimensions change.
\setlength{\emergencystretch}{2em}

\begin{document}

\title[DMGH: Decoupling Mutation Publication from Graph Maintenance]{DMGH: Decoupling Mutation Publication from Graph Maintenance for Dynamic HNSW Serving}

\author{Ding Qidong}
\email{dingqidong@stu.xjtu.edu.cn}

\author{Zhang Shuai}
\email{s.z@stu.xjtu.edu.cn}
\affiliation{%
  \institution{Department of Computer Science, Xi'an Jiaotong University}
  \city{Xi'an}
  \country{China}
}

\author{Wu Weiguo}
\email{wgwu@xjtu.edu.cn}

\author{Nie Shiqiang}
\email{shiqiang.nie@xjtu.edu.cn}
\affiliation{%
  \institution{Department of Computer Science, Xi'an Jiaotong University}
  \city{Xi'an}
  \country{China}
}

\renewcommand{\shortauthors}{Ding Qidong et al.}

\begin{abstract}
Continuous insertions, deletions, and vector updates require dynamic HNSW to maintain retrieval quality and low query tail latency while incorporating object mutations. We present DMGH, a decoupled execution framework for dynamic HNSW serving. DMGH separates mutation publication from graph maintenance, provides supplementary search entry points for objects to which the graph has not yet adapted, and processes query candidates using the object states and vectors read by the query. In the background, local repair and batched integration adapt search paths, switch between old and new structures, and reclaim old representations while queries continue. This design allows insertions, deletions, and updates to publish their respective object changes, reducing the dependence of foreground operations on complete structural maintenance.

All 24 test cells comprising six million-scale and ten-million-scale dataset configurations and four workloads achieve checkpoint Recall@10 $\geq0.95$ and a median query P99 across repetitions of $\leq20$ ms. In the comparison on update-heavy workloads over three million-scale datasets, DMGH is the only system among the tested configurations to meet both targets on all datasets, with a completed mutation rate of 5.35K--10.05K/s. At the two operating points where both DMGH and CleanANN meet the quality threshold, DMGH improves query throughput by 3\%--9\% and reduces P99 by 27\%--34\%. In three-hour mixed-workload runs, DMGH improves query throughput by 7.9\%--19.3\% over the configuration with lifecycle maintenance disabled, while continuing graph integration and snapshot reclamation.

\end{abstract}

\keywords{dynamic approximate nearest-neighbor search; HNSW; mutation publication; local graph repair; safe memory reclamation}

\maketitle



\section{Introduction}
\label{sec:introduction}

Approximate nearest-neighbor (ANN) search is a fundamental operation in recommendation, dense retrieval, and retrieval-augmented generation~\cite{covington2016youtube,karpukhin2020dpr,lewis2020rag}. HNSW provides efficient navigation through a hierarchical proximity graph and has become an important index structure for vector search~\cite{malkov2020hnsw,pan2024vdbsurvey}. Objects in online services are continuously inserted and deleted, and the vector of an existing object is updated as its content changes. The index must incorporate these mutations while maintaining effective search paths and low tail latency~\cite{mohoney2025quake,wang2026candor,dean2013tail}.

All three mutation types involve object state and navigation structure: after publishing a new object, insertion still requires establishing search entry points; after revoking result eligibility, deletion still requires cleaning up connections and storage; after replacing a vector, an update still requires adaptation to its new location. Performing all of this work synchronously makes the foreground wait for neighbor search, connection adjustment, and storage reorganization; deferring structural maintenance requires specifying how queries handle changing objects during the transition. Existing dynamic ANN methods incorporate object mutations through local graph updates, maintenance during access, or background consolidation~\cite{zhang2026cleanann,yu2025greator,singh2021freshdiskann}. For HNSW serving under continuous insertions, deletions, and updates, separating publication from maintenance leaves three questions: how candidates are discovered and processed before graph adaptation is complete, how subsequent maintenance adjusts search paths, and how new structures replace old representations while queries continue.

A same-label update illustrates the distinction between object readability and reachability through graph paths: after a vector moves from $x_i$ to $x^{\prime}_i$, a query can read the new vector once it accesses the object, but adjacency organized around the old location may prevent the object from being discovered. Section~\ref{sec:publication-repair} uses this scenario to diagnose path mismatch; new objects likewise require discovery entry points before graph integration, while deleted objects must be filtered by state before structural cleanup.

DMGH advances mutation publication, transitional candidate processing, and subsequent graph maintenance separately. Queries discover candidates through committed dynamic slots and hints associated with existing graph nodes, filter deleted objects according to the states read, and handle vector replacement through update-overlay rules. The background reserves capacity for repair connections within bounded adjacency lists, adapts local paths using HNSW neighbor pruning, and integrates dynamic objects in batches over the committed prefix. The graph representation is published before integration progress is advanced; existing queries retain old representations, and physical reclamation waits for the relevant readers to exit.

The experiments examine serving performance, structural effects, and maintenance costs separately. All 24 cells formed by six 1M or 10M dataset configurations and four workloads meet checkpoint R95 and a median query P99 of at most 20 ms. On SIFT1M and DEEP1M under update-heavy workloads, both DMGH and CleanANN meet R95; DMGH achieves 1.03--1.09 times the QPS and 27\%--34\% lower P99. Paired experiments with bypass candidates disabled further demonstrate the effect of local adjacency repair on graph paths to moved objects; three-hour mixed-workload comparisons quantify the effects of lifecycle processing on serving performance and the accumulation of historical state.

The main contributions are as follows:
\begin{enumerate}
\item \textbf{Publication and candidate processing under continuous insertions, deletions, and updates.} Dynamic entry points supplement object discovery before graph integration, while state observation and vector reads handle deletions and updates. Within the candidate budget, explicit rules replace distances for updated objects and remove old base results that have not been rescored, allowing query processing to proceed without waiting for complete graph maintenance.
\item \textbf{Degree-constrained local repair and batched graph integration.} Complete adjacency construction combines reserved repair edges with HNSW neighbor pruning. Ordered publication of the graph root and prefix cursor transfers committed dynamic slots into graph navigation in batches while retaining old representations accessed by existing queries.
\item \textbf{Quantification of serving performance and maintenance costs across mutation workloads.} Dynamic serving is evaluated under insert-heavy, delete-heavy, update-heavy, and mixed workloads. Targeted audits, paired graph-traversal experiments, module ablations, and three-hour runs examine object processing, path recovery, and the benefits and costs of maintenance.
\end{enumerate}

\section{Background and Motivation}
\label{sec:background}

\subsection{Graph-Based ANN Search}
\label{sec:graph-ann}

Given an object set $\mathcal{O}=\{o_1,\ldots,o_n\}$, each object is associated with a vector $x_i\in\mathbb{R}^d$. For a query $q$ and distance function $D$, ANN returns an object set approximating the true top-$K$ at lower computational cost. Let $R_K(q)$ be the set of valid labels in the first $K$ positions of the returned list, and let $T_K(q)$ be the true nearest-neighbor label set. Then
\begin{equation}
\label{eq:recall}
\operatorname{Recall@}K=\frac{|R_K(q)\cap T_K(q)|}{K}.
\end{equation}
HNSW first locates a region through sparse upper layers, then maintains a candidate set and expands neighbors at the bottom layer~\cite{malkov2020hnsw}. Search width affects recall and the costs of distance computation and adjacency access~\cite{aumuller2020annbench,azizi2025grapheval}. Object state and search paths serve different roles in a graph index: object state determines eligibility for inclusion in query results, whereas graph connections affect whether a query can discover relevant objects within a finite search budget. Thus, publishing a new or updated vector does not imply that the corresponding search paths have been established, and making a deletion mark effective does not imply that the associated graph connections have been cleaned up.

\subsection{Dynamic Object Semantics}
\label{sec:dynamic-semantics}

Dynamic search must distinguish logical objects from physical representations. A logical object is identified by a stable label \(i\), which remains unchanged by vector updates and identifies the object in query results. Physical representations include the storage instance containing the vector and its structural position in the index. Because update publication, structural adjustment, and reclamation of old storage can proceed separately, one logical object may temporarily have multiple physical representations. The Recall definition above takes the set of labels appearing in the returned positions, so each label contributes at most one hit.

An object's logical state is determined by its stable label, liveness state, and current vector when live. The liveness states are absent, live, and deleted. We use mutation as a collective term for three operations: \texttt{insert(i,x)} creates an object and its initial vector; \texttt{delete(i)} revokes its eligibility for inclusion in results; and \texttt{update(i,x')} replaces its vector while retaining its label.

Object mutations involve two related but distinct processes. State publication makes a new liveness state or vector the basis for subsequent reads; structural adaptation adjusts candidate entry points and adjacency relationships to the changed object set and vector distribution. For example, deletion may first revoke result eligibility and then clean up the corresponding structures; a vector update may first publish a new payload and then adjust adjacency associated with its spatial location; a newly inserted object may first participate in queries through a temporary entry point and then be integrated into the long-term graph structure. This distinction means that state publication determines how queries interpret objects already discovered, whereas structural adaptation affects whether queries can discover those objects effectively. The two processes may complete at different times but must be coordinated during dynamic serving.

\subsection{Object Observation and Serving Objectives}
\label{sec:consistency-state}

DMGH reads object state and payload on a per-candidate basis. A query first checks the label's liveness state and then obtains a payload pointer. Candidates observed as deleted are filtered out; distances for the remaining candidates are computed after a nonnull payload is obtained. The state check determines candidate eligibility, whereas the payload read determines the input to the current scoring operation; each occurs at its own observation point. Each distance computation uses the same immutable payload instance that was obtained.

An update occurring after a query obtains a payload affects subsequent reads by replacing the pointer; existing readers continue using the instance they have obtained. Adjacency and graph-root replacement use the same access protection: new structures become available to subsequent reads, while old structures are not freed until the relevant readers exit~\cite{brown2015debra}. Section~\ref{sec:authoritative-resolution} gives the object-processing order, and Section~\ref{sec:unified-lifecycle} describes the retirement and reclamation paths for different representations.

Retrieval quality is evaluated at checkpoint $b$. The reference set consists of the live labels and their payloads at that checkpoint, from which the true top-$K$ is computed. Query throughput and tail latency characterize serving during execution, while the completed mutation rate characterizes foreground progress. Counts of integrated structures, pending work, and retired and reclaimed objects characterize background progress. Section~\ref{sec:evaluation-methodology} specifies the aggregation rules for these metrics.

\subsection{Motivation}

\subsubsection{Serving Stalls from Synchronous Graph Maintenance}
\label{sec:sync-stalls}

Graph updates typically involve neighbor search, edge selection, adjacency modification, and concurrency synchronization. If these operations share a large critical section with queries, a single update may block other requests; if full rebuilding is performed in a concentrated phase, maintenance costs may be incurred over a short interval. Both organizations can make structural maintenance directly affect online serving.

Figure~\ref{fig:motivation-native-paths} compares two synchronous organizations: graph updates enclosed by a coarse-grained lock (global-lock wrapper) and periodic rebuilding (periodic rebuild). All four configurations achieve Recall@10 of 0.979--0.988, but the global-lock insert and mixed configurations obtain only approximately 0.8 QPS and approximately 10 s P99; its update configuration achieves a mutation throughput of 0.090/s. The periodic-rebuild update configuration exhibits 110.86 s P99.9 and a mutation throughput of 2.536/s. These blocking diagnostics demonstrate the importance of execution organization: even an index with high recall may struggle to serve continuously under synchronous maintenance.

\begin{figure}[t]
  \centering
  \includegraphics[width=0.96\columnwidth]{figures/figure1.pdf}
\caption{Online serving performance under two structural-maintenance organizations. GL and PR denote the global-lock wrapper and periodic rebuild, respectively; \textbf{(a)} query throughput, \textbf{(b)} query tail latency P99/P99.9, and \textbf{(c)} completed mutation throughput. All configurations achieve Recall@10 \(\geq 0.95\); the vertical axes use logarithmic scales.}
\Description{Comparison of query throughput, tail latency, and completed mutation throughput between the global-lock wrapper and periodic rebuild under the same experimental conditions.}
  \label{fig:motivation-native-paths}
\end{figure}

\subsubsection{Semantic Visibility and Structural Discoverability}
\label{sec:publication-repair}

Once structural maintenance is removed from state publication, object-state changes and search-path adjustments can occur at different times. Insertions and updates must address candidate discovery during this interval, whereas deletions affect the result eligibility of discovered objects through state checks. We examine the former using a same-label vector update: when the new vector has been published, adjacency may still correspond to the pre-update location. Even though a query can read the new vector after accessing the object, it may fail to reach the object because the paths are no longer suitable.

To examine this distinction, we move a set of objects to a substantially different vector region and retrieve them solely through graph traversal, holding the initial graph, query set, and search budget fixed. Across three random seeds, Recall@1 for the moved objects drops from 1.0 before the update to 0 after publication of the new payloads.

This result shows that, in this vector-migration scenario, publishing a new payload changes the distance computed after the object is accessed but does not restore graph paths to its new location. Separating publication from maintenance therefore requires transitional discovery entry points for inserted or updated objects, followed by graph adaptation. Section~\ref{sec:evaluation-local-repair} examines path recovery through local repair; the targeted audit in Section~\ref{sec:semantic-audit} checks deletion filtering and the returned states after publication of the three mutation types.

\subsection{Design Requirements}
\label{sec:design-principles}

These observations yield three requirements. First, transitional retrieval after mutations must be supported: new objects need supplementary entry points before graph integration, deleted objects are filtered according to the state observed by the query, updated objects are scored using the vectors read, and the handling of old results under insufficient budgets must be explicit. Second, navigation must be adapted incrementally: local repair allocates space to existing neighbors and repair connections within node degree limits, while dynamic objects are integrated in batches over the committed prefix. Finally, structural switching must be separated from physical deallocation: queries access new structures according to published progress, and old representations held by existing readers are retained until their accesses finish.

\begin{figure*}[t]
  \centering
  \includegraphics[width=0.88\textwidth]{figures/figure2.pdf}
\caption{DMGH's coordinated execution pipeline. The foreground makes new object representations available to queries, while the background connects structural-adaptation results to the same query path. Stable logical labels link the stages, and reader epochs protect physical accesses during handoffs between old and new representations.}
\Description{DMGH architecture: the foreground publishes object states, queries access objects through multiple physical entry points, and the background maintains adjacency, promotes dynamic objects, and manages reclamation of old representations.}
  \label{fig:dmgh-overview}
\end{figure*}

\section{DMGH Design}
\label{sec:overview}

DMGH is implemented in C++ on top of Faiss-HNSW~\cite{douze2024faiss}, reusing hierarchical traversal and distance-computation primitives. The system targets continuous insertions, deletions, and updates, and distinguishes foreground-published object changes from subsequently maintained navigation structures. Figure~\ref{fig:dmgh-overview} shows the shared state accessed by the foreground, queries, and background. This section defines the publication unit of each operation; the next two sections detail transitional candidate processing and structural maintenance, respectively.

\subsection{State Organization}
\label{sec:state-organization}

DMGH manages object state and navigation structures separately, allowing object mutations to be published before graph adaptation. The base graph \(B\) provides the underlying search paths, while the state table \(S\) and update-payload table \(U\) store the states and updated payload pointers of base and promoted objects, respectively. Object-state or vector changes therefore need not wait for adjacency adjustment to finish. Newly inserted objects are first stored in dynamic slots \(X\); once initialized, they become accessible to queries through the commit boundary.

While navigation structures have not yet adapted, anchor-associated candidate hints \(C\) provide supplementary discovery entry points. The background publishes local repair results through complete adjacency snapshots \(D\) and integrates committed dynamic slots into the graph in batches. This latter process is termed promotion, and the resulting promoted graph is managed by \(P\): the graph root provides a search entry point, and the prefix cursor marks the range of slots already integrated into the graph. Replaced representations and their retirement epochs are recorded in \(R\) and reclaimed after the relevant readers exit.

Different entry points associate the same object through its stable logical label. Queries read state and payload from \(S/U\) or \(X\) according to the region containing the label. A dynamic slot stores the initial payload, subsequent update pointer, object state, and initial neighbors, and is allocated at a monotonically increasing offset. The logical label consists of the base\_label fixed at the first insertion plus the slot offset; deletion, promotion, and reclamation never reuse that offset. A payload version denotes an immutable payload instance produced by an insertion or update. Publishing a new instance changes the target of subsequent payload reads, while existing readers continue using the instance already acquired within their protected access interval.

\subsection{Execution Paths and Handoffs}
\label{sec:end-to-end-ops}

\textbf{Foreground.} Insert initializes $X$ and advances the commit boundary; update replaces $U[i]$ or the payload pointer of a dynamic slot; delete writes the object state. The publication point $t_{\mathrm{pub}}$ and API return point $t_{\mathrm{return}}$ mark the representation switch and call completion, respectively. Each branch also performs the work listed in Table~\ref{tab:foreground-background-contract} before returning.

\textbf{Queries.} A query accesses $B/D$ under the protection of sharded read views and a reader epoch, and merges candidates provided by dynamic slots, $P$, and $C$. New objects participate through committed entry points. Candidates obtain states and vectors through $S/U/X$, and objects observed as deleted are filtered out. The update overlay replaces or removes existing distances for the same label, after which results are sorted and written out.

\textbf{Background.} Local repair discovers anchors and publishes candidate hints. When adjacency materialization and complete-snapshot mode are enabled, it selects neighbors under the degree constraint and replaces complete adjacency snapshots in $D$. Versioned promotion constructs $P$ in batches and publishes the graph root before the prefix cursor. Retirement records $R$ manage old storage after structural replacement. Section~\ref{sec:repair-lifecycle} details these two structural-maintenance paths.

\subsection{Publication and Reclamation Discipline}
\label{sec:foreground-background-contract}

Each representation is published only after construction is complete. Payload pointers and commit boundaries allow subsequent reads to obtain complete objects, while adjacency-snapshot and graph-root replacements allow subsequent traversals to obtain new structures. Retirement records retain replaced storage until the relevant readers exit. Table~\ref{tab:foreground-background-contract} distinguishes state publication from API return and specifies the work that remains on the foreground after publication.

\begin{table}[t]
\centering
\caption{Publication and return boundaries of foreground operations. The background continues unfinished structural integration and storage reorganization.}
\label{tab:foreground-background-contract}
\begingroup\footnotesize
\begin{tabularx}{\columnwidth}{@{}p{0.19\columnwidth}XX@{}}
\toprule
Operation & Published object & Work before return \\
\midrule
Insert & Commit boundary of initialized slots & Wait for gates; construct anchors; submit hints and reciprocal work after publication \\
Dynamic Update & Slot update pointer & Check state; retire old instances and process hints \\
Base / promoted Update & $U[i]$ payload pointer & Check state; account for old payloads and submit or execute repair \\
Delete & Slot state or $S[i]$ & Branch-specific sampled repair, routing-state handling, and payload cleanup \\
\bottomrule
\end{tabularx}\endgroup
\end{table}

The next section details the three mutation types and candidate processing; Section~\ref{sec:repair-lifecycle} presents the rules for structure selection, publication, and reclamation of old representations.

\section{Mutation Publication and Candidate Processing}
\label{sec:foreground-resolution}

Algorithm~\ref{alg:foreground-mutation} specifies mutation publication order, and Algorithm~\ref{alg:search-merge-filter} specifies candidate merging and update-overlay processing within the budgets.

\subsection{Foreground Mutation Publication}
\label{sec:foreground-publication}

The evaluation uses one mutation worker; queries, promotion, and maintenance may overlap with foreground execution. Dynamic-label indices in Algorithm~\ref{alg:foreground-mutation} refer to the corresponding slots, and insertion returns base\_label plus the slot offset.

\begin{algorithm}[t]
\caption{ForegroundMutation}
\label{alg:foreground-mutation}
\KwIn{$op$; optional label $i$, vector $x$; one mutation worker}
\uIf{$op=\AState{INSERT}$}{
    Wait for base-mutation and promotion gates\;
    $j\gets\AFn{FetchAdd}(X.next\_slot,1)$; build payload and anchors\;
    Initialize $X[j]$ with payload, anchors, LIVE state, null update\;
    $\AFn{ReleaseFetchMax}(X.committed\_boundary,j+1)$\;
    Publish hints; enqueue reciprocal work; \Return{$X.base\_label+j$}\;
}
\uElseIf{$op=\AState{UPDATE}$}{
    Enter base-reader epoch and mutation gates\;
    \eIf{$i$ is unpromoted}{
        Require LIVE; build $p$; read old slot update pointer\;
        $\AFn{ReleaseStore}(X.update[i],p)$; retire old standalone payload\;
        Publish hints or promoted anchors\;
    }{
        \If{$S[i]$ is deleted or redirected}{Leave scope; \Return{}\;}
        Append $p$; $p_{old}\gets\AFn{Exchange}(U[i],p)$\;
        Account $p_{old}$; enqueue or run repair\;
    }
    Leave mutation scope\;
}
\ElseIf{$op=\AState{DELETE}$}{
    \If{$i$ is unpromoted}{Store $X.state[i]\gets\AState{DELETED}$; \Return{}\;}
    Run sampled repair if selected\;
    Store TOMBSTONE to $S[i]$ and routing state\;
    Exchange $U[i]$ to null; retire old; clear repair snapshot\;
}
\end{algorithm}

\textbf{Anchor construction and branch differences.} Insert searches the base graph for initial anchors, retaining valid, nondeleted, distinct base labels in their original ranking, up to the smaller of the neighbor budget and the base size. An empty set is returned if there is no vector, no base object, or a zero budget. Update selects the slot pointer or $U[i]$ according to whether the object has been promoted, retiring a standalone instance or accounting for the old payload in append-only storage, respectively. The base branch of Delete may perform sampled repair before publishing a tombstone and clearing the payload and repair snapshot. The internal redirect branch separately writes the object and routing redirection states and clears the old payload and snapshot.

\subsection{State Observation and Payload Scoring}
\label{sec:authoritative-resolution}

Candidates produced by insertions, deleted labels remaining in the structure, and objects with updated vectors all locate their states and payloads through logical labels: base and promoted objects use $S/U$, while unpromoted objects use $X$. When processing an initial candidate $i$, a query first checks its deletion state, then obtains a nonnull payload and computes the distance. Let $t_i^s$ and $t_i^p$ denote the state-read and payload-read times, respectively. A candidate is processed as follows:
\begin{equation}
\label{eq:resolve}
\begin{aligned}
 z_i&\gets\operatorname{ReadState}(i;t_i^s),\\
 p_i&\gets\operatorname{ReadPayload}(i;t_i^p),\quad z_i\ne\mathrm{deleted},\\
 s_i(q)&=\operatorname{dist}(q,\operatorname{vec}(p_i)),\quad p_i\ne\varnothing.
\end{aligned}
\end{equation}
Payload reads use acquire semantics, with $t_i^s<t_i^p$. A candidate is skipped if deleted is observed. If deletion occurs after the state check, the payload is still read: a null pointer causes the candidate to be skipped, whereas a nonnull instance is used for scoring. A new instance published after $p_i$ is acquired does not change the input to this distance computation. The reader epoch covers payload use, sorting, and result output; release of old instances waits for the relevant readers to exit.

\subsection{Multi-source Candidate Discovery and Merge}
\label{sec:multi-source}

A query first traverses the base graph and active adjacency, then incorporates the promoted graph, dynamic-slot hints, repair hints, and the base and promoted update overlays in that order, and finally scans the remaining dynamic slots. Each entry point locates state and payload by label. Hints select objects for evaluation, overlays handle old distances, and fallback scanning supplements slots not evaluated through dynamic hints.

\begin{algorithm}[t]
\caption{SearchMergeFilter}
\label{alg:search-merge-filter}
\KwIn{Query $q$, size $k$, budgets $b$}
\KwOut{Top-$k$ entries; state/payload reads follow Eq.~\ref{eq:resolve}}
Wait for base-mutation/compaction gates; enter reader views and epoch\;
$\mathcal{C}_q\gets\AFn{TraverseBaseAndAdjacency}(q,b.traversal)$\;
$\mathcal{R}_q\gets\AFn{FilterAndScore}(\mathcal{C}_q,q,b.rerank)$\;
Acquire committed boundary $c$ and cursor $p$, then acquire and pin root $g$\;
Append enabled traversal results from $g$\;
$J\gets\AFn{ProbeDynamicHints}(q,p,c,b.hint)$; $V\gets\varnothing$\;
\ForEach{$j\in J$ in ascending slot order}{
    Check state, score available payload, append; record evaluated slot in $V$\;
}
$H\gets\AFn{ProbeRepairHints}(q,b.repair)$; remove labels in $\mathcal{R}_q$\;
Evaluate eligible $H$ within exact budget; quality-order and filter duplicates\;
Append repair candidates within acceptance budget\;
\ForEach{updated label $i$ within each region scan budget (base, then promoted)}{
    Remove all entries labeled $i$; check state and append available payload score\;
}
Remove old base results with an update overlay not covered by the scan\;
\ForEach{slot $j\in[p,c)$ within the position-counted fallback budget}{
    \If{$j\notin V$}{Check state, score available payload, append\;}
}
Partially sort by (distance, label); write first $k$ entries\;
Release root, epoch and reader views\;
\end{algorithm}

\textbf{Update-overlay processing.} When a base or promoted label is scanned, all entries with that label are first removed from the result list $\mathcal{R}_q$; the state is then checked, and the acquired payload is rescored. For a base label with an update overlay that is not covered by this scan, its old base result is removed. Thus, a limited scan budget does not fill results by retaining old base distances. $\mathcal{R}_q$ denotes the list of result entries, distinct from the shared hints $C$ and retirement records $R$.

\textbf{Anchor-associated hints.} To locate candidates within a bounded number of probes, queries use anchors to select hint buckets. Query anchors preferentially use explicit inputs; otherwise, they are drawn from current search results and may be expanded along adjacency. In scored mode, scores are initialized with routing distance plus a dimension-dependent far-distance penalty; penalties are accumulated for repeated anchors, and the list is truncated in ascending score and label order. Unscored mode preserves first-occurrence order. This procedure is distinct from the base-graph anchor search used by insertion and background repair.

Dynamic-slot hints map anchors to 4096 buckets. Bucket IDs are deduplicated during publication, and a full bucket cyclically overwrites its oldest record. Queries select buckets under the anchor-count budget, deduplicate bucket IDs, and visit them in ascending order, reading records from newest to oldest within each bucket. Slots beyond the commit boundary are excluded; in versioned-promotion mode, encountering the promoted prefix ends the scan of that bucket. Slots are deduplicated across buckets, probing stops when the budget is reached, and selected slots are scored in ascending slot-ID order. Repair hints use separate first-in, first-out records. With score ordering enabled, records are read in descending hint-score order, descending epoch order, and ascending label order; otherwise, newer records are read first. Labels are deduplicated across buckets. Hint scores determine candidate priority, while final result distances are obtained by scoring payloads.

\textbf{Duplicate handling.} Dynamic-slot hints deduplicate by slot, and the fallback scan skips slots already evaluated through these hints. Repair hints filter existing labels before evaluation and during quality ordering, and update overlays replace entries with the same label. Final sorting retains the list of result entries without global cross-source label deduplication or minimum-distance merging. The old-cursor/new-root combination described in Section~\ref{sec:repair-to-promotion} allows the same label to enter the top-$k$ through both graph and dynamic entry points; such duplicate entries may occupy multiple returned positions.

\subsection{Query Budgets and Result Assembly}
\label{sec:query-budgets}

Separate budgets constrain candidate discovery, slot scanning, distance computation, and result acceptance. efSearch controls graph search width, while the rerank budget limits scoring of initial candidates. Dynamic-hint probing is bounded by the number of slots deduplicated across buckets, and repair-hint probing by the number of unique labels. Fallback and other exact scans count visited slot positions; anchor and repair-candidate scoring count actual distance computations; the repair acceptance budget counts labels added after quality ordering. Filtered slots consume scan budget but do not necessarily consume scoring budget. The anchor limit counts input labels, with zero indicating no limit; bucket capacity bounds the number of retained records.

\section{Graph Reintegration and Safe Reclamation}
\label{sec:repair-lifecycle}

Local repair replaces adjacency per node, prefix promotion integrates dynamic objects in batches, and lifecycle management subsequently reorganizes payloads and reclaims old representations.

\subsection{Repair Inputs and Replacement Construction}
\label{sec:repair-obligation}

\textbf{Repair inputs.} Each step dequeues requests within its budget, skips deleted triggering objects, and searches for anchors in batches using current payloads. On failure, it falls back to individual searches, with validity and deduplication rules as in Section~\ref{sec:foreground-publication}. Anchors are used to publish hints and, when materialization is enabled, to generate adjacency replacements. The triggering object is the event owner, while replacement targets may include multiple nodes.

\textbf{Degree-constrained adjacency selection.} Repair connections compete with existing neighbors for limited adjacency space. Let $m$ be the degree limit for the current construction at target node $u$. Construction first removes invalid labels, self-loops, out-of-range entries, and duplicates from the repair neighbors, and excludes deleted dynamic labels, yielding $O_u$. These neighbors are then merged with valid, distinct bottom-layer neighbors of the base graph; exact symmetric distances are computed, and HNSW neighbor pruning produces a candidate list $H_u$. To admit prioritized repair connections into the final adjacency, DMGH reserves capacity for $O_u$ before admitting $H_u$.

The reservation is determined by the node's adjacency deficit and repair activity. Let $a_u$ denote the degree-deficit bonus, $q_u$ the heat bonus, and $z_u$ an optional hit bonus. Then
\begin{equation}
\label{eq:repair-reserve}
r_u=\min\!\left(\rho_0+a_u+q_u+z_u,\rho_{\max},|O_u|,m\right).
\end{equation}
$\rho_0$ and $\rho_{\max}$ are the base reservation and its upper bound, respectively. $a_u$ is the floor of the nonnegative difference between the base graph's bottom-layer adjacency capacity and its current degree, divided by $s_d$. Each incremental edge-patch publication increments raw heat by one and refreshes its timestamp. At read time, heat is reduced by the integer part of the elapsed logical-clock count divided by the decay step, with a floor of zero; dividing the decayed heat by $s_h$ and rounding down yields $q_u$. Only when both the hit penalty and materialization bonus are enabled is $z_u$ the integer part of the decayed hit bonus divided by $s_h$; otherwise, it is zero. If $s_d$ or $s_h$ is zero, its corresponding bonus is zero; if $O_u$ is empty or $m=0$, then $r_u=0$.

The experiments use $\rho_0=1$, $\rho_{\max}=3$, $s_d=4$, and $s_h=2$, with the hit bonus disabled. Complete snapshots can invoke adjacency construction directly; the reservation is still computed from the decayed heat of the actual patch-publication history.

Repair neighbors are ordered by descending materialization bonus, ascending distance, and ascending label; with the bonus disabled in the experiments, they are ordered by distance and label. Construction admits the first $r_u$ repair neighbors, previously unselected neighbors from $H_u$, and then the remaining repair neighbors, stopping at $m$ to produce a complete replacement adjacency list.

\textbf{Repeated requests and queue capacity.} Repair requests for the same object are coalesced while queued or being processed to reduce redundant work. With priorities enabled, a queued request retains the larger displacement and refreshes its sequence number; otherwise, the original request is retained. Requests arriving during processing are buffered as one record, combining the larger displacement and latest sequence number, after which re-enqueueing is attempted when the current processing finishes. Requests for new objects are dropped when the queue is full; buffered requests are also dropped if re-enqueueing fails, and the object's registration is released. Neither case falls back to synchronous repair. Request-processing completion, hint publication, and adjacency publication are recorded as separate progress events.

\subsection{Adjacency-snapshot Replacement}
\label{sec:adjacency-delta}

With materialization and complete snapshots enabled, the last replacement for each target is retained, and all immutable blocks are allocated before any exchange. If the target is observed to be deleted before publication, its block is discarded. Otherwise, $D[u]$ is exchanged under the shard lock, the old block is retired with the current epoch, and the lock is released. If deletion is observed after publication, the active snapshot is cleared. Targets are published separately, and existing readers continue accessing old blocks. Incremental mode appends edge patches; hints can still be published with materialization disabled. Algorithm~\ref{alg:maintenance-lifecycle} follows the complete-snapshot path.

\subsection{Committed-prefix Promotion}
\label{sec:repair-to-promotion}

Prefix promotion publishes the graph root before advancing the cursor, preventing queries from skipping slots before the new graph is available. A publication attempt must acquire the mutex and close the dynamic-write gate. If either step fails, pending state is retained and the attempt ends. On success, it waits for registered dynamic writers to finish, then reads the commit upper bound $c$; ordinary queries need not exit.

Each round processes $[p,e)$ within the slot budget, where $p$ is the current cursor and $e\leq c$. Nothing is published unless the minimum batch size is met, and deleted slots also count toward the budget. After construction, vector and node storage, anchor entry points, and $P.root$ are published in that order, followed by a release store advancing $P.cursor$ to $e$. Each maintenance step attempts promotion only once; if the entire step makes no progress, the current wake-up ends and the work is retried later.

A query first reads the commit boundary and cursor with acquire semantics, then acquires the graph root and retains a shared reference to it. The graph corresponding to a new cursor has already been published. An old cursor may also be paired with a new graph root, allowing graph and slot entry points to overlap; duplicates are handled as described in Section~\ref{sec:multi-source}. Shared references and epochs protect representations held by existing queries.

A repair request completes when its processing finishes; hint and adjacency publications are recorded separately. Promotion completes when the graph root has been published and the cursor advanced. Append-to-base promotion uses different gating; the algorithms here detail the versioned path.

\subsection{Lifecycle Phases and Reclamation}
\label{sec:unified-lifecycle}

Algorithm~\ref{alg:maintenance-lifecycle} advances repair, promotion, and storage reorganization in sequence, and also consumes deferred reciprocal-connection work within a budget; processing is deferred during a pressure-induced pause. Standalone payloads and old adjacency blocks are reclaimed when their retirement epoch is less than the safe epoch, while old graph roots additionally wait for shared references to be released. Superseded payloads in append-only storage are first accounted for and later reclaimed through physical compaction. The experiments report integration, compaction, reclamation, and resource costs separately.

\begin{algorithm}[t]
\caption{MaintenanceLifecycleStep}
\label{alg:maintenance-lifecycle}
\KwIn{Budgets $b$; full snapshots and versioned promotion enabled}
\ForEach{$e\in\AFn{DequeueAndCoalesce}(b.repair)$}{
    \If{owner observed deleted}{Consume $e$; \Continue\;}
    Search anchors from current payload; publish repair hints\;
    Build complete blocks by the reserve-and-prune rule (\S\ref{sec:repair-obligation})\;
    Keep last replacement per target; allocate all blocks\;
    \ForEach{replacement $(u,A_u)$}{
        \If{$u$ observed deleted}{Discard $A_u$; \Continue\;}
        Lock shard; $A_{old}\gets\AFn{Exchange}(D[u],A_u)$\;
        Retire $A_{old}$ with current epoch; unlock shard\;
        \If{$u$ now observed deleted}{Clear its active repair snapshot\;}
    }
    If an in-flight update was recorded, retry enqueue; on failure drop and unregister\;
}
\If{promotion pending and publish mutex/gate acquired}{
    Wait for registered dynamic writers; read cursor $p$ and committed $c$\;
    \If{minimum publication batch is met}{
        Select $e\leq c$ by slot budget; build $g_{new}$ for $[p,e)$\;
        Publish vector/node storage and anchor heads\;
        $g_{old}\gets\AFn{Exchange}(P.root,g_{new})$\;
        $\AFn{ReleaseStore}(P.cursor,e)$; retire $g_{old}$ with epoch\;
    }
    Open dynamic gate; release mutex\;
}
Retain pending promotion on failed acquisition\;
Advance budgeted reciprocal work and physical compaction\;
Reclaim epoch-safe objects subject to retained root references; emit telemetry\;
\end{algorithm}

\section{Evaluation}
\label{sec:evaluation}

The evaluation examines overall serving, concurrent object processing, search costs, path recovery, and online integration in sequence, and quantifies maintenance benefits and costs through ablations and three-hour runs.

\subsection{Experimental Setup and Methodology}
\label{sec:evaluation-methodology}

\textbf{Platform and implementation.} All experiments run on a single-socket Intel Xeon 6505P with 12 cores, 24 hardware threads, and approximately 32 GiB DRAM, using Ubuntu 22.04 and Linux 6.8.0-111-generic. DMGH is implemented on Faiss-HNSW. Short-duration main experiments run for 30 s with 3 independent repetitions, using 8 query threads and 1 mutation thread. Background workers run according to each system's configuration and are included in the CPU and memory costs of the same process. Unless otherwise stated, all systems run on the same machine using their native memory or storage layouts.

\textbf{Datasets, queries, and workloads.} The 1M experiments use SIFT1M, DEEP1M, and GIST1M, and the 10M experiments use SIFT10M, DEEP10M, and SPACEV10M. Each repetition starts from a frozen index containing only base vectors, with the same initial logical live set $L(0)$ across systems. Query threads repeatedly execute fixed test queries during the timed window. Mutation streams are divided into four workloads by operation mix: insert-heavy consists entirely of insertions; delete-heavy consists of 5\% insertions, 90\% deletions, and 5\% updates; update-heavy consists of 5\% insertions, 5\% deletions, and 90\% updates; and mixed consists of 34\% insertions, 33\% deletions, and 33\% updates. Query and mutation rates are computed from their respective completed-operation counts.

\textbf{Quality checkpoints and statistics.} At evaluation time $t_e$, the authoritative live set $L(t_e)$ is frozen, deleted labels are excluded, and the corresponding payloads of updated objects and committed insertions are used to compute exact L2 top-10 results for 100 fixed queries. Recall@10 $\geq 0.95$ is abbreviated as R95. For three independent repetitions $r=1,2,3$, the joint target R95-QDT@20ms is defined as
\begin{equation}
\label{eq:quality-target}
\min_r\operatorname{Recall@10}_r\geq0.95,\qquad
\operatorname{median}_r(\mathrm{P99}_r)\leq20\,\mathrm{ms}.
\end{equation}
QPS, P99, P99.9, and completed Mutation/s are reported as medians across repetitions. The three-hour experiments summarize serving metrics over a single complete run and report endpoint Recall separately.

\textbf{Comparison protocol.} We compare against CleanANN, Greator, LSM-VEC, and CONDA, retaining each system's native storage organization, threading model, and update API. Each system's configuration is selected using a separate tuning set and its native parameter space, after which the parameters are frozen. These experiments compare end-to-end serving performance under the native configurations; query-performance ratios are computed only at operating points where both systems meet R95. Completed Mutation/s is reported separately as the observed rate within the 30 s window. Systems complete different numbers of mutations, and checkpoint Recall corresponds to each system's own dynamic live set. Faiss global-lock, periodic rebuild, and lightweight publication paths provide comparisons of execution organizations of the same type.

\subsection{End-to-End Dynamic Serving}
\label{sec:end-to-end}

\subsubsection{Dynamic Serving Across Dataset--Workload Cells}
\label{sec:matrix}

We first examine whether the same dynamic serving path covers different data distributions, dataset sizes, and mutation types. Figure~\ref{fig:eval-dynamic-matrix} summarizes the dynamic serving matrix formed by six datasets and four workloads. All 12 cells at 1M meet R95-QDT@20ms: SIFT1M achieves 2.3K--4.5K QPS and 2.0--4.6 ms P99, and DEEP1M achieves 2.0K--4.7K QPS and 3.0--5.6 ms P99; the high-dimensional GIST1M still sustains 1.9K--2.4K QPS and 7.4--8.3 ms P99. All 12 cells at 10M also pass: SIFT10M and DEEP10M achieve approximately 1.8K--3.4K QPS and 3.4--5.7 ms P99, while SPACEV10M achieves approximately 1.7K QPS and 15.0 ms P99 with a larger search budget. These results show that the unified execution path accommodates joint variation in scale, dimensionality, and mutation composition.

\begin{figure*}[t]
\centering\setlength{\pairedfigureheight}{0.27883\textwidth}
\begin{minipage}[t]{0.56122\textwidth}
\vspace{0pt}\centering
\includegraphics[height=\pairedfigureheight]{figures/figure721.pdf}
\caption{DMGH's 24-cell dynamic serving matrix. (a)--(d) show 1M and (e)--(h) show 10M; panels report repeat-min Recall@10, repeat-median QPS, P99, and Mutation/s in order. Dashed lines indicate Recall@10 $=0.95$ and P99 $=20$ ms.}
  \label{fig:eval-dynamic-matrix}
\Description{Recall, query throughput, P99, and mutation throughput for six datasets and four workloads.}
\end{minipage}\hfill
\begin{minipage}[t]{0.41878\textwidth}
\vspace{0pt}\centering
\includegraphics[height=\pairedfigureheight]{figures/figure722.pdf}
\caption{External comparison on three 1M update-heavy datasets. (a) Recall@10; (b) QPS; (c) P99; (d) completed Mutation/s over 30 s. Gray cross-hatched bars identify operating points that do not meet R95.}
  \label{fig:eval-external}
\Description{End-to-end serving comparison of DMGH, CleanANN, Greator, LSM-VEC, and CONDA on three datasets.}
\end{minipage}
\end{figure*}



\subsubsection{Comparison with CleanANN and Other External Systems}
\label{sec:external}

Figure~\ref{fig:eval-external} compares three 1M update-heavy operating points. DMGH is the only system among the tested configurations to meet both R95 and P99 $\leq20$ ms at all three points, with Recall@10 of 0.955--0.987, P99 of 4.568--8.279 ms, and a completed mutation rate of 5,353.7--10,050.7/s over 30 s.

On SIFT1M and DEEP1M, where both systems meet R95, DMGH improves QPS over CleanANN by 3\% and 9\%, respectively; P99 decreases from 6.943 and 7.698 ms to 4.568 and 5.630 ms, reductions of 34.2\% and 26.9\%. CleanANN achieves Recall of 0.891 on GIST1M, so this point is excluded from performance ratios under the shared quality threshold. Greator meets R95 at all three points, but its P99 is 24.681--62.666 ms. LSM-VEC meets both targets only on SIFT1M. CONDA achieves P99 of 0.686--5.492 ms, but its Recall of 0.770--0.885 falls below R95 at all three points.

\subsection{Foreground Semantics and Quality--Cost Trade-offs}
\label{sec:semantics-quality}

\subsubsection{Concurrent Publication Semantics}
\label{sec:semantic-audit}

The concurrent audit checks object discovery after insertion, target filtering after deletion, and vector identification after updates. It starts with 4,096 base vectors and performs three repetitions, each with 100 mutation rounds and 4 continuously running query threads, using $k=10$ and efSearch of 64. Each round invokes insertion, update, and deletion in sequence; after all calls return, a timestamp is recorded and the round number is published. After acquiring this number, query threads search for the corresponding labels using the inserted vector, the deleted target's original base vector, and the updated vector, respectively; these three searches form one scored query group. Target labels are disjoint across rounds and operation types, and later rounds do not change the expected states of earlier targets. Each round waits for at least 100 query groups or until the 2,000 ms deadline.

An insertion hit and a stale return after deletion are determined by whether the top-$k$ contains the respective target label. The updated target's current vector equals the query, while its old base vector differs by approximately 5,000 in each dimension. Returned entries for the target label with distances below $10^{-4}$ are therefore classified as current-version entries, and the remaining entries as old-version entries. These two events are counted independently, and both are counted if they occur in the same query. Each of the four rates uses the number of query groups exhibiting the corresponding event as its numerator and the total number of scored query groups $N$ as its denominator. Let $D$, $O$, $T$, and $R$ denote the numbers of query groups with stale returns after deletion, query groups with old-version returns after updates, timed-out rounds, and mutation rounds, respectively. The aggregate violation rate is
\begin{equation}
\label{eq:audit-violation}
V=\frac{D+O+T}{2N+R}.
\end{equation}
Individual insertion misses are not separately included in the numerator; rounds in which no jointly correct observation of all three mutations has been recorded by the end of the wait contribute to $T$. Round-level visibility latency starts at the recorded time after all three calls in the round have returned and ends when the first query group completes with an insertion hit, no return of the deleted target, a current-version update hit, and no old-version update return. It includes that group's query execution time. Quantiles are computed only over successful rounds, and timed-out rounds are reported separately.

Each of the three repetitions checks at least 10,416 query groups, with active query groups sampled before mutations in at least 99 rounds. In every repetition, the insertion hit rate and update current-version hit rate are both 1.0, while the deletion stale-return rate, update old-version return rate, and aggregate violation rate are all 0. All 100 rounds achieve a jointly correct observation, with no timeouts. The median of the per-repetition P99 values of round-level visibility latency is 1,137.481 $\mu$s, and the maximum is 1,203.649 $\mu$s.

\subsubsection{Search-Budget Trade-offs}
\label{sec:quality-cost}

\begin{figure*}[t]
\centering\setlength{\pairedfigureheight}{0.26751\textwidth}
\begin{minipage}[t]{0.61922\textwidth}
\vspace{0pt}\centering
\includegraphics[height=\pairedfigureheight]{figures/figure732.pdf}
  \caption{Search-budget trade-offs on SIFT1M and SPACEV10M: (a) QPS; (b) P99; (c) Mutation/s. Filled and hollow markers denote repeat-min and repeat-median Recall@10, respectively.}
  \label{fig:eval-quality-cost}
\Description{Retrieval quality and serving costs under different search budgets on SIFT1M and SPACEV10M.}
\end{minipage}\hfill
\begin{minipage}[t]{0.36078\textwidth}
\vspace{0pt}\centering
\includegraphics[height=\pairedfigureheight]{figures/figure751.pdf}
\caption{Execution-path ablation under R95. \textbf{(a)} repeat-min Recall@10; \textbf{(b)} QPS; \textbf{(c)} P99/P99.9; \textbf{(d)} Mutation/s. The dashed line indicates R95, and throughput and latency plots use logarithmic axes.}
  \label{fig:eval-path-ablation}
\Description{Recall, query throughput, tail latency, and mutation throughput for six execution-path configurations.}
\end{minipage}
\end{figure*}

Figure~\ref{fig:eval-quality-cost} jointly varies efSearch/rerank budget. On SIFT1M, 32/32 achieves the highest QPS but only 0.898 repeat-min Recall. At 48/48, median Recall is 0.953 but repeat-min Recall is 0.948, still below R95. The first passing configuration is 64/64, with repeat-min Recall of 0.962 and 2,677.0 QPS; 96/96 raises Recall to 0.982 while reducing QPS by 14.6\%.

On SPACEV10M, repeat-min Recall is 0.901 at 96/64 and 0.939 at 256/256; 384/384 is the first configuration to reach 0.951. Raising the budget to 512/512 increases Recall to 0.962 but reduces QPS by 26.4\%. The budget required to meet the same quality target varies by dataset, and additional search work also reduces throughput.

\subsection{Local Repair and Online Graph Reintegration}
\label{sec:structural-convergence}

\subsubsection{Repair-Induced Reachability Recovery}
\label{sec:evaluation-local-repair}

This section uses same-label vector migration as a diagnostic scenario to test whether degree-constrained adjacency replacement can restore graph paths to objects. Three paired structural-damage experiments run repair-on and repair-off using the same initial graph, moved objects, and query set within each pair. Bypass candidate publication is disabled, so recovery can occur only through the adjacency structures actually published.

\begin{table}[t]
  \centering
\caption{Path recovery through local repair and online integration of dynamic objects. (a) Three paired repair experiments, with each pair sharing the initial graph, random seed, and query set; (b) maintenance endpoints and cumulative counts from three online runs.}\label{tab:eval-repair}
  \scriptsize\setlength{\tabcolsep}{3pt}
\textbf{(a) Paired repair diagnostics}\par\smallskip
  \resizebox{\columnwidth}{!}{%
  \begin{tabular}{@{}lrrr@{}}
    \toprule
    Audit metric & Repair on & Repair off & Across-pair summary \\
    \midrule
    Final Recall@1 & $\geq$ 0.969 & 0 & +0.985 \\
    Graph-reachable moved-object fraction & 1.000 & 0 & --- \\
    Structural-region fraction & --- & --- & +0.996 \\
    Recovery latency & 390.7 ms & N/A & --- \\
    Final QPS & 2,933.2 & 5,531.1 & --- \\
    Final P99 & 0.383 ms & 0.208 ms & +0.175 ms \\
    Candidate publications & 0 & 0 & 0 \\
    Repair queue pending at termination & 0 & 0 & 0 \\
    \bottomrule
  \end{tabular}}
  \par\medskip
\textbf{(b) Online integration statistics}\par\smallskip
  \begin{tabular*}{\columnwidth}{@{\extracolsep{\fill}}lrrr@{}}
    \toprule
Metric & \shortstack{DEEP10M\\update-heavy} & \shortstack{DEEP10M\\mixed} & \shortstack{SPACEV10M\\mixed} \\
    \midrule
Duration (s) & 600 & 300 & 300 \\
Final promotion backlog & 0 & 5 & 1 \\
Cumulative promoted objects & 304,981 & 142,578 & 29,248 \\
Cumulative graph probes & N/A & 635,854 & 204,824 \\
    \bottomrule
  \end{tabular*}
  \par\smallskip
  \begin{minipage}{\columnwidth}
\scriptsize In (b), repair pending, inflight, and debt are all 0 at the endpoints of all three runs. Graph probes denote promoted-graph probes; N/A indicates that this metric was not collected in that run.
  \end{minipage}
\end{table}

Table~\ref{tab:eval-repair}(a) shows that repair-on restores the graph-reachable fraction to 1 for all three seeds, with Recall@1 of at least 0.969; both metrics are 0 with repair-off. The median paired gain in structural-region coverage is 0.996, and the median recovery latency is 390.664 ms. Both configurations have 0 candidate publications, so recovery arises from the published adjacency paths. Repair increases search work: QPS decreases from 5,531.050 to 2,933.222, and P99 increases from 0.208 to 0.383 ms. This experiment measures path recovery and its cost; full serving performance is reported in Sections~\ref{sec:end-to-end} and~\ref{sec:architectural-ablation}.

\subsubsection{Online Reintegration of Repaired and Promoted States}
\label{sec:promoted-traversal}

Table~\ref{tab:eval-repair}(b) reports three online-integration runs: DEEP10M update-heavy for 600 s, and DEEP10M and SPACEV10M mixed for 300 s each, sampled every 5 s. The final promotion backlog is at most 5 in all runs, and repair pending, inflight, and debt are all 0. The two mixed runs record 635,854 and 204,824 promoted-graph probes, respectively. Promotion counts reflect integration progress, while probe counts show that published entry points participate in queries.

\subsection{Architecture and Module Ablations}
\label{sec:architectural-ablation}

\subsubsection{Publication Cost and Maintenance Trade-offs}
\label{sec:path-ablation}



Figure~\ref{fig:eval-path-ablation} compares six within-system operating points on SIFT1M update-heavy, all meeting R95. Global-lock and periodic rebuild achieve 0.266 and 0.149 QPS, respectively, with P99 of 30.030 s and 53.867 s, exhibiting severe blocking under synchronous maintenance. Full DMGH achieves 3,750.5 QPS and 3.119 ms P99.

Minimal overlay-only achieves 4,595.3 QPS and 2.192 ms P99; the corresponding values for minimal + routing are 4,137.5 QPS and 2.401 ms. Full executes repair, promoted traversal, and lifecycle maintenance together. Relative to the former, its QPS is 18.4\% lower and P99 is 42.3\% higher, quantifying the short-run query cost of complete maintenance.

Relative to w/o lifecycle, Full improves QPS and Mutation/s by 15.4\% and 18.4\%, respectively, and reduces P99 by 13.6\%, but increases P99.9 from 5.469 to 6.856 ms. Maintenance benefits must therefore be evaluated separately for throughput, the ordinary tail, and the extreme tail, in conjunction with the workload ablations and long-run experiments below.

\subsubsection{Contributions of Individual Modules}
\label{sec:module-ablation}

% Ablation verification pending: the user specifies the promoted-graph ablation as "disabling query-side traversal of the promoted graph."
% Complete switch combinations for w/o lifecycle and Foreground-only still require the existing run configurations; do not infer the disabled scope from configuration names.
Module ablations use efSearch/rerank = 96/96, 8 query threads, and a 30 s, 3-repetition protocol across four SIFT1M workloads. All 24 configurations meet R95, with quality margins of 2.3--3.9 percentage points.

Removing lifecycle maintenance reduces QPS by 18.9\%--38.7\%, increases P99 by 19.7\%--41.2\%, and reduces Mutation/s by 0.9\%--8.3\% across the four workloads; on delete-heavy, P99.9 increases from 5.239 ms to 16.430 ms. These comparisons show benefits from lifecycle maintenance for throughput and P99 across all four workloads, while changes in extreme tail latency depend on the workload.

Repair and promoted traversal act on local structures and query entry points, respectively, and their benefits vary with workload composition. The query-side traversal switch controls probes of the promoted graph; the background may still construct and publish roots. Removing local repair reduces QPS by 10.4\% on insert-heavy and 12.4\% on delete-heavy; on mixed, QPS changes by less than 1\%, while P99.9 increases by 48.8\%. Disabling query-side traversal of the promoted graph reduces QPS by 9.0\% on insert-heavy and 12.2\% on delete-heavy; P99 increases by 12.3\% on update-heavy, and P99.9 increases by 45.5\% on mixed.

Foreground-only reduces QPS by 17.6\%--38.4\% and increases P99 by 7.4\%--44.9\%; under insert-heavy, Mutation/s decreases by 69.0\%, although some configurations have lower P99.9. Complete maintenance improves throughput and P99 in this comparison, while its effect on the extreme tail varies by workload.

\begin{figure*}[t]
  \centering
  \includegraphics[width=\textwidth]{figures/figure752.pdf}
\caption{Module ablations on SIFT1M. \textbf{(a)} quality margin above R95; \textbf{(b)--(e)} changes in QPS, P99, P99.9, and Mutation/s relative to Full DMGH under the same workload. All cells meet R95; Recall is reported as repeat-min and the remaining metrics as repeat-median.}
\Description{Quality margins and throughput and latency changes relative to Full for six module configurations under four workloads.}
  \label{fig:eval-module-ablation}
\end{figure*}

\subsection{Long-run Serving and Lifecycle Progress}
\label{sec:longrun}

This section runs Full DMGH and w/o lifecycle for 3~h each under mixed workloads on SIFT1M, DEEP1M, and GIST1M to examine serving performance and background progress under continuous mutations. Each configuration executes one complete run. Serving and maintenance states are sampled every 30~s, and process resources approximately every 5~s. Figure~\ref{fig:longrun-timeseries} shows the time series, Figure~\ref{fig:longrun-performance} reports aggregate performance and endpoint Recall for these runs, and Figure~\ref{fig:longrun-endpoint} reports endpoint states and resource costs.

\subsubsection{Long-run Service Stability}
\label{sec:longrun-service}

Figure~\ref{fig:longrun-timeseries}(a)--(c) shows that Full DMGH sustains relatively high query and mutation throughput on all three datasets. Comparing the means of the first and last 60 sampling windows (approximately 30~min each), its QPS changes by less than 1\% on every dataset, and P99 exhibits no clear sustained increase. In contrast, w/o lifecycle on SIFT1M decreases from approximately 6.81K to 6.51K QPS, while P99 increases from 1.439~ms to 1.497~ms; its query throughput on GIST1M remains largely stable.

Figure~\ref{fig:longrun-performance} shows that Full improves QPS by 7.9\%--19.3\% and Mutation/s by 8.7\%--17.6\%, while reducing P99 and P99.9 by 6.7\%--12.6\% and 4.3\%--13.3\%, respectively. In the order SIFT1M, DEEP1M, and GIST1M, Full achieves endpoint Recall of 0.994, 0.958, and 0.977, compared with 0.974, 0.955, and 0.969 for the ablation; all meet R95. The time series describe serving during the run, whereas Recall is evaluated on the endpoint dynamic live set.
\subsubsection{Deferred-State Absorption and Reclamation}
\label{sec:longrun-lifecycle}

Figure~\ref{fig:longrun-timeseries}(e)--(f) shows the progress of structural integration and compaction during the runs. The three Full DMGH runs promote 9.42M, 6.96M, and 1.78M dynamic objects and complete 2,141, 1,580, and 403 compactions. Both counts are zero for w/o lifecycle. Promotion counts integrated objects, whereas compaction counts commit events. Window rates are computed by dividing differences in cumulative counts by the actual window duration, reflecting the progress of integration and compaction, respectively.

In Figure~\ref{fig:longrun-endpoint}(a)--(b), Full has endpoint promotion backlogs of 7, 6, and 2, compared with approximately 8.02M, 6.40M, and 1.53M for the ablation. Full reclaims 548,094, 404,274, and 102,632 repair snapshots cumulatively, leaving only 2, 0, and 0 retired snapshots at the endpoint. The ablation records no reclamation and retains 466,100, 371,995, and 88,006 retired snapshots, respectively. Completed work and remaining state jointly reflect the background's continuing processing of deferred state.

Figure~\ref{fig:longrun-endpoint}(c)--(f) quantifies the resource costs of complete maintenance. Full has higher endpoint dynamic payload and memory amplification than the ablation, while process Peak RSS is lower by 1.4\%, 21.4\%, and 11.3\%, respectively. In particular, DEEP1M decreases from 23.79 GiB to 18.69 GiB and GIST1M from 28.97 GiB to 25.68 GiB. The two configurations complete different numbers of mutations. Dynamic payload describes payload storage, snapshot counts describe the disposition of structural representations, and Peak RSS summarizes resident memory for the entire process. RSS traces use the process timeline, including construction and final quality evaluation; serving and maintenance rates use the online-workload timeline. The endpoint values in Figure~\ref{fig:longrun-endpoint}(c)--(d) are also the sampled peaks in these runs; (e) reports the process peak from the run summary, distinct from the maximum discrete sample in Figure~\ref{fig:longrun-timeseries}(d). Maintenance CPU is computed as CPU time divided by the 3~h duration, expressed as utilization of one core.

Full's maintenance CPU corresponds to 6.37\%, 6.09\%, and 4.90\% of one core, compared with 7.30\%, 6.28\%, and 4.44\% for the ablation. Full completes more mutations and retains more payload storage, while leaving fewer objects pending promotion and fewer retired snapshots, with lower peak RSS. These metrics describe work volume, remaining state, and resource costs, respectively.

\begin{figure*}[t]
  \centering
  \includegraphics[width=\textwidth]{figures/figure761.pdf}
  \caption{Serving and maintenance time series for 3~h mixed workloads on three datasets: (a) QPS; (b) P99; (c) Mutation/s; (d) RSS; (e) Promotion/s; (f) Compaction/s. Colors identify datasets; solid/dashed lines denote Full/w/o lifecycle. Serving and maintenance use approximately 30~s windows; RSS is sampled approximately every 5~s. Timeline definitions are given in the text.}
  \label{fig:longrun-timeseries}
\Description{Three-hour time series of queries, mutations, resources, and maintenance on three datasets.}
\end{figure*}

\begin{figure*}[t]
  \centering\setlength{\pairedfigureheight}{0.30824\textwidth}
  \begin{minipage}[t]{0.45472\textwidth}
    \vspace{0pt}\centering
    \includegraphics[height=\pairedfigureheight]{figures/figure762.pdf}
    
    \captionof{figure}{3~h mixed-run summaries: (a) endpoint Recall@10 (dashed: R95); (b) QPS; (c) P99/P99.9; (d) Mutation/s. Each configuration has one run; Recall uses the endpoint live set.}
    \label{fig:longrun-performance}
  \end{minipage}\hfill
  \begin{minipage}[t]{0.52528\textwidth}
    \vspace{0pt}\centering
    \includegraphics[height=\pairedfigureheight]{figures/figure763.pdf}
    
    \captionof{figure}{3~h endpoint states and costs: (a) promotion backlog (symlog); (b) live/reclaimed/retired repair snapshots; (c) dynamic payload; (d) memory amplification; (e) process Peak RSS; (f) maintenance CPU as utilization of one core. Snapshot counts are normalized by the sum of the three categories, with totals above the bars; live/retired are endpoint counts, and reclaimed is cumulative.}
    \label{fig:longrun-endpoint}
  \end{minipage}
\Description{Serving summaries, endpoint quality, structural states, and resource costs for three-hour runs.}
\end{figure*}

\section{Related Work}
\label{sec:related-work}

\subsection{Graph-Based Approximate Nearest-Neighbor Search}
\label{sec:related-graph}

Graph-based ANN indexes organize candidate navigation through proximity connections. NN-Descent constructs approximate nearest-neighbor graphs using neighborhood propagation and iterative refinement~\cite{dong2011nndescent}; HNSW provides navigation through hierarchical small-world graphs~\cite{malkov2014nsw,malkov2020hnsw}; NSG optimizes connectivity, paths, and out-degree~\cite{fu2019nsg}; and ParlayANN studies parallel graph construction~\cite{manohar2024parlayann}. DiskANN, SPANN, and Starling scale retrieval through aspects such as disk layout, partitioning, and block access~\cite{subramanya2019diskann,chen2021spann,wang2024starling}. DMGH reuses HNSW navigation and neighbor pruning to study candidate processing, local structural adaptation, and dynamic-object integration under continuous insertions, deletions, and updates.

\subsection{Dynamic and Streaming ANN Systems}
\label{sec:related-dynamic}

CleanANN adapts to distribution changes through workload-aware connections, handles deletions through query-adaptive neighborhood consolidation, and uses lock-free, semi-lazy cleanup~\cite{zhang2026cleanann}. DMGH separates foreground publication from graph maintenance, handles old distances through budgeted update-overlay processing, and combines background repair with graph integration to support continuous insertions, deletions, and updates.

IP-DiskANN handles insertions and deletions in place~\cite{xu2025ipdiskann}; Wolverine repairs path damage caused by deletions~\cite{liu2025wolverine}; Greator limits the scope of connection and disk-page modifications~\cite{yu2025greator}; and CONDA combines connectivity-aware pruning with lazy deletion~\cite{lee2026conda}. DMGH combines bounded repair reservations with HNSW pruning and publishes complete adjacency per node.

FreshDiskANN incorporates recent mutations through temporary indexes, deletion lists, query merging, and StreamingMerge~\cite{singh2021freshdiskann}; LSM-VEC combines neighbor search, connection writes, and LSM buffering and compaction~\cite{zhong2025lsmvec,oneil1996lsm}. DMGH retains the HNSW base, uses committed slot prefixes as the unit of integration, and coordinates graph and dynamic entry points by publishing the graph root before the cursor.

The partition-based systems SPFresh and Quake accommodate updates through local rebalancing and adjustable partitions, respectively~\cite{xu2023spfresh,mohoney2025quake}. FGIM studies efficient merging of multiple existing graph indexes~\cite{wu2026fgim}, while CANDOR-Bench systematically examines continuous ingestion, distribution drift, and concurrent access~\cite{wang2026candor}. These works and DMGH advance dynamic vector serving through index organization, merging operations, and evaluation methodology.

\subsection{Concurrent Representation Management}
\label{sec:related-lifecycle}

The Bw-tree uses indirect mapping and consolidation to manage incremental representations~\cite{levandoski2013bwtree}. Hazard pointers, epoch-based reclamation, and DEBRA protect access to old objects~\cite{michael2004hazard,brown2015debra}; batched deallocation and allocator interactions can also affect tail latency~\cite{kim2024batchfree}. DMGH uses per-node adjacency exchanges and batched graph-root publication as structural switching points, with epochs and shared root references constraining the release of old representations.

\section{Conclusion}
\label{sec:conclusion}

DMGH separates object publication from graph maintenance under continuous insertions, deletions, and updates, connecting object changes to navigation adaptation through transitional candidate processing, degree-constrained local repair, and ordered prefix promotion. All 24 serving cells meet R95-QDT@20ms. In external comparisons where both systems meet R95, DMGH improves throughput over CleanANN by 3\%--9\% and reduces P99 by 27\%--34\%. Paired experiments demonstrate the path-recovery effect of local repair, while three-hour mixed-workload comparisons show that lifecycle maintenance improves throughput by 7.9\%--19.3\% and sustains graph integration and reclamation of old representations.

\FloatBarrier
\label{page:body-end}

\bibliographystyle{ACM-Reference-Format}
\bibliography{references}

\end{document}
